\section{NAMED ENTITY RECOGNITION \& COLLOCATION ANALYSIS}

\subsection{Problem Definition}
Before retrieved textbook passages can improve LLM accuracy, the FAISS retrieval query must be precise. Raw clinical vignettes dilute the embedding with incidental language. NER-based query rewriting isolates only the clinically meaningful tokens before querying FAISS.

\subsection{Model Choice}
The en\_ner\_bc5cdr\_md model \cite{neumann2019} is trained on the BC5CDR corpus \cite{li2016}, a collection of 1,500 PubMed articles annotated for 5,818 disease and 4,409 chemical entity types. It produces two typed entity labels: DISEASE and CHEMICAL, mapping directly to USMLE test goals.

\subsection{Corpus Analysis}
Running en\_ner\_bc5cdr\_md over all 10,178 MedQA training questions produced 54,256 total entities (DISEASE: 39,575 / CHEMICAL: 14,681), with a mean of 5.33 entities per question.

\begin{figure}[htbp]
    \centering
    \includegraphics[width=0.7\textwidth]{entities.png}
    \caption{Entity type frequencies extracted by en\_ner\_bc5cdr\_md.}
    \label{fig:entity_freq}
\end{figure}

\begin{figure}[htbp]
    \centering
    \includegraphics[width=\textwidth]{entity_dist.png}
    \caption{Left: distribution of entity counts. Right: top 20 most frequently extracted clinical entities.}
    \label{fig:entity_dist}
\end{figure}

\subsection{Retrieval Score Impact}
FAISS top-1 retrieval scores were compared for raw vignettes versus NER-rewritten entity strings across all three embedding models. NER rewriting helps both biomedical-scale models: Qwen3-Embedding-0.6B (+0.005) and Octen-Embedding-0.6B (+0.006), but hurts the general-purpose MiniLM model (-0.022).

\begin{figure}[htbp]
    \centering
    \includegraphics[width=\textwidth]{ner_impact.png}
    \caption{Mean top-1 FAISS retrieval score for raw vignettes vs. NER-rewritten entity strings.}
    \label{fig:ner_impact}
\end{figure}

\subsection{Collocation Analysis}
PMI and chi-square bigram analysis was applied to all MedQA training question text. A PMI-derived lay-to-clinical thesaurus was constructed by filtering bigrams where one word appears in a curated lay-term seed list.

\begin{figure}[htbp]
    \centering
    \includegraphics[width=\textwidth]{bigrams.png}
    \caption{Top 25 clinical bigrams by PMI score and chi-square score.}
    \label{fig:bigrams}
\end{figure}